% ===========================================================================
% Section 3: Data and evaluation protocol; Section 4: Method
% ===========================================================================

\section{Data and Evaluation Protocol}\label{sec:data}

\subsection{Split and Domains}\label{sec:split}

The original BDD100K train folder has 69,863 images with detection labels. They were split once, with seed 20260916, into 55,000 train images and 14,863 held-out test images, stratified jointly by time of day, weather and scene type. The official validation set (val, 10,000 images) is kept unchanged. The public 20,000-image test set has no detection labels and is therefore not used. The membership of each split is frozen by manifests and SHA256 hashes. No image is shared between splits, either by file name or by content.

The source domain is daytime with clear weather. The three target domains form a $2\times2$ design over time of day and weather (Table~\ref{tab:domains}): lighting changes only the illumination, weather changes only the weather, and compound changes both.

\begin{table}[t]
\centering
\caption{Domains and split sizes.}
\label{tab:domains}
\footnotesize
\setlength{\tabcolsep}{3pt}
\begin{tabular}{@{}l>{\raggedright\arraybackslash}p{0.52in}r>{\raggedright\arraybackslash}p{0.75in}>{\raggedright\arraybackslash}p{1.0in}@{}}
\toprule
Domain & Condition & \begin{tabular}[t]{@{}r@{}}Train\\pool\end{tabular} & Final test & Selection val \\
\midrule
source   & daytime, clear & 9,804  & 2,650 & 1,764 (source checkpoint selection) \\
lighting & night, clear   & 18,015 & 4,869 & 3,274 (schedule pilot only) \\
weather  & daytime, rainy & 1,985  & 933 (537 test + 396 val) & none \\
compound & night, rainy   & 1,738  & 756 (470 test + 286 val) & none \\
\bottomrule
\end{tabular}
\end{table}

The val images of weather and compound were moved into the final test set, because the original test sets of these two domains (537 and 470 images) were below the minimum of 600 images. As a consequence, weather and compound have \textbf{no val set for selection}; every target-domain run keeps the final update of a fixed schedule instead of selecting a checkpoint.

\subsection{Labels}\label{sec:labels}

The labels are the 2018 BDD100K detection release, obtained through the \texttt{solesensei\_bdd100k} mirror on Kaggle; every image and label file matches the mirror in size and CRC32. The 2018 release has been superseded by the Detection 2020 labels (\texttt{det\_20}). Keeping the 2018 release is a recorded decision, and it is stated among the limitations (Section~\ref{sec:limitations}). The legacy class names (person, bike, motor) are mapped to the canonical names.

\subsection{Zero-Shot Gate}\label{sec:gate}

Before any ladder was run, the zero-shot gap of each shift was measured with the source model: source val versus the target domain's train pool, which the model had never seen, so no test label was read. A shift whose 95\% CI does not exceed 0.01 AP is flagged. All three shifts pass the gate (Table~\ref{tab:gate}).

\begin{table}[t]
\centering
\caption{Zero-shot gate (source val versus the target train pool, 1,000 bootstrap resamples).}
\label{tab:gate}
\footnotesize
\setlength{\tabcolsep}{3pt}
\begin{tabular}{@{}lcc@{}}
\toprule
Shift & Zero-shot gap (AP) & 95\% CI \\
\midrule
lighting & 0.100 & 0.091--0.110 \\
weather  & 0.032 & 0.024--0.041 \\
compound & 0.129 & 0.121--0.139 \\
\bottomrule
\end{tabular}
\end{table}

Weather passes the gate, but with the smallest gap. Every weather result is therefore reported with its CI and the raw change in AP.

\subsection{Test-Label Access}\label{sec:testaccess}

No run reads test labels during training; the manifest of each run records this. After a group of runs has finished, all of the group's prediction files are frozen in a plan (with the hash of every file and every checkpoint), and that plan is then scored exactly once. There are three plans (Table~\ref{tab:plans}).

\begin{table}[t]
\centering
\caption{The three frozen plans, each scored exactly once.}
\label{tab:plans}
\footnotesize
\setlength{\tabcolsep}{3pt}
\begin{tabular}{@{}l>{\raggedright\arraybackslash}p{1.25in}cl@{}}
\toprule
Plan & Contents & \begin{tabular}[t]{@{}c@{}}Prediction\\sets\end{tabular} & Scored on \\
\midrule
campaign      & fine-tuning, 3-epoch ceiling, COCO control & 87 & 6 Nov.~2026 \\
RQ5           & \aldi and replay under the original rules & 46 & 7 Nov.~2026 \\
supplementary & 12-epoch ceiling, additional replay seeds, matched-passes experiment & 42 & 23 Nov.~2026 \\
\bottomrule
\end{tabular}
\end{table}

All three plans use the same 1,000 test-image resamples (seed 20261030), so every comparison between arms is paired on the same resamples.

\section{Method}\label{sec:method}

\subsection{Detector and Source Model}\label{sec:detector}

The detector is Faster R-CNN with a ResNet-50 FPN backbone (Lin et al.~\cite{lin2017fpn}; the torchvision implementation), initialized from COCO (Lin et al.~\cite{lin2014coco}) and run at the native resolution of $1280\times720$ (\texttt{min\_size=720}, \texttt{max\_size=1280}). This architecture was chosen because of RQ5: almost every public DAOD implementation targets Faster R-CNN, so choosing it keeps a comparison with an adaptation method feasible.

The source model is trained once on the 9,804 source images: 12 epochs, batch 4, automatic mixed precision (AMP), learning rate (LR) 0.005, a warmup of 500 updates, 29,412 updates in total (about 118,000 image passes). The checkpoint is selected on source val only; the selected checkpoint is the final update, with a reliable mAP$_{50:95}$ of 0.352 on source val and 0.342 on source test. Every target-domain arm is initialized from this same checkpoint.

\subsection{Fine-Tuning Schedule}\label{sec:schedule}

Every target-domain fine-tuning run (the budget rungs, the ceilings, the controls) uses one shared rule: train for $\max(S \text{ updates}, E \text{ epochs})$ on the $N$ sampled images, then keep the final update. The fixed settings are batch 2, SGD with momentum 0.9, weight decay $5\cdot10^{-4}$, LR $10^{-3}$, a warmup of 50 updates followed by cosine decay to 1\% of the LR, gradient clipping at 10, and AMP.

The values of $S$ and $E$ were fixed by a pilot on lighting val only, under a rule written before the pilot ran: take the largest $S$ and the smallest $E$ that lie within $\delta = 0.005$ AP of the best configuration. The result is \textbf{\boldmath $S = 500$, $E = 3$}. The value $E = 3$ was chosen by a very thin margin: $E = 2$ trails $E = 5$ by exactly 0.00511 AP, that is, it lies outside $\delta$ by only 0.0001. Consequences of the rule:
\begin{itemize}
\item $N \leq 333$ receives 500 updates; for $N = 10$ this is equivalent to about 100 epochs.
\item $N \geq 500$ receives $1.5\cdot N$ updates (3 epochs at batch 2).
\item The pilot already indicated that at $N = 10$, fine-tuning makes the model 0.05--0.06 AP worse than the source model. This is recorded as a result, not as a reason to retune the schedule.
\end{itemize}

\subsection{Budget Ladder, Ceilings and Controls}\label{sec:ladder}

\begin{itemize}
\item \textbf{Ladder:} $N = 10$, 25, 50, 100, 250, 500, 1,000 and the full pool; lighting additionally has 2,000, 5,000 and 10,000. Every rung below the full pool has 3 sampling seeds (0, 1, 2), nested by seed: for the same seed, the $N = 25$ sample contains the $N = 10$ sample.
\item \textbf{3-epoch ceiling:} fine-tuning on the full pool under exactly the schedule above, 2 seeds. This is the top rung of the ladder.
\item \textbf{12-epoch ceiling} (added in response to the review): full pool, 12 epochs (the same number of epochs as the source model), 2 seeds, every other setting unchanged. On lighting it is 108,096 updates, about 216,000 image passes, compared with about 54,000 for the 3-epoch ceiling.
\item \textbf{Matched pool of 1,738 images} (for RQ4): the lighting and weather ceilings are recomputed at the pool size of compound, for both the 3-epoch and the 12-epoch ceiling, 2 seeds.
\item \textbf{COCO-initialization control:} at $N = 100$ and $N = 1{,}000$, the same run but initialized from the COCO weights instead of the source model, seed 0.
\end{itemize}

\subsection{Metric, Recovery and Thresholds}\label{sec:metrics}

\textbf{The primary metric} is mAP$_{50:95}$ averaged over a class roster that is fixed for every domain: pedestrian, car, traffic light, traffic sign (called \emph{reliable mAP}). These are the four classes with the most instances in all four domains; the metric is undefined if a class is missing. The 10-class mean is a sensitivity analysis. Every score goes through a single evaluation path (pycocotools).

\textbf{Drop} is defined as
\begin{equation}
\mathrm{Drop} = \mAPsrc - \mAPzero,
\label{eq:drop}
\end{equation}
where \mAPsrc is the score of the source model on source test and \mAPzero is the score of the source model on the target domain's test set (zero-shot).

\textbf{Recovery(\boldmath$N$)} is defined as
\begin{equation}
\mathrm{Recovery}(N) = \frac{\mathrm{mAP}(N) - \mAPzero}{\mAPc - \mAPzero},
\label{eq:recovery}
\end{equation}
where \mAPc is the target-domain ceiling, not \mAPsrc. A target domain can be intrinsically harder than the source domain (night is harder than day), so normalizing by \mAPsrc would mix the domain gap with intrinsic difficulty. Recovery is not clipped: a negative value means that fine-tuning is worse than zero-shot. The minimum denominator is 0.01 AP; when more than 5\% of the bootstrap resamples have a denominator below this level, recovery and the thresholds are recorded as \emph{undetermined}, together with the number of such resamples.

\textbf{Thresholds} are defined in a sustained sense on the recovery curve averaged over seeds:
\begin{itemize}
\item \NBE: the smallest rung $N > 0$ at and above which recovery is always $> 0$ (break-even).
\item \Nsf, \Nnf: the smallest rung $N$ at and above which recovery is always $\geq 0.75$ or $\geq 0.90$, respectively.
\item A threshold that is reached only at the ceiling rung itself (where recovery equals 1 by definition) is recorded as \textbf{not determined within the pool}, together with the largest rung that has not reached it. The $N_{95}$ threshold of our original study proposal has been dropped, because it falls into exactly this case in all three shifts.
\end{itemize}

\textbf{Uncertainty} is estimated with a hierarchical bootstrap: test images are resampled (1,000 times, seed 20261030) and sampling seeds are resampled; the seed-to-seed variation of the ceiling is also propagated into the denominator. The interval of each threshold is taken from the 95\% CI band of the recovery curve.

\textbf{The RQ1 interaction} is $\mathrm{Drop}(\text{compound}) - \mathrm{Drop}(\text{lighting}) - \mathrm{Drop}(\text{weather})$. Besides the CI, we report the minimum detectable interaction, $(z_{0.975} + z_{0.80})\cdot\mathrm{SE}$ (SE is the bootstrap standard deviation), both in AP and as a fraction of the weather effect.

\textbf{The log-linear fit},
\begin{equation}
\mathrm{Recovery} = \alpha + \beta\cdot\ln N,
\label{eq:fit}
\end{equation}
is used only for interpolation. The fit region runs from the first rung whose lower CI bound is above 0 (and stays so at every larger rung) to the largest rung below the ceiling; at least 4 rungs are required. $R^2$ and RMSE are computed on the fit region only. A threshold from the fit is reported only when it lies within the corresponding ladder interval. Held-out check: the fit is made on the three lowest rungs of the region and predicts the remaining rungs.

\subsection{\texorpdfstring{\aldi}{ALDI++} and the Replay Control Arm}\label{sec:aldi}

\aldi was reimplemented in torchvision from the paper and the published configuration files (the original code repository has no license, so its code was not copied). The configuration follows \aldi: an EMA teacher; labeled images receive strong augmentation and random erasing; on unlabeled images, the student receives strong augmentation plus MIC-style block masking (ratio 0.5, 32 px blocks), and the teacher receives a weakly augmented version; soft distillation on all four outputs (classification and regression of the ROI head, objectness and regression of the RPN); hard pseudo-labels (threshold 0.8) used only for anchor and proposal sampling; no feature alignment; the scored model is the EMA teacher.

\textbf{Differences from the original \aldi:} fixed $1280\times720$ images (no multi-scale training); no separate burn-in stage, because every arm is initialized from the same source model as FT; a small batch, 2 labeled and 2 unlabeled images per step, with LR $10^{-3}$ as for FT; the use of labeled target images (semi-supervised) is our extension, since the original ALDI uses no target labels.

\begin{table}[t]
\centering
\caption{The arms of RQ5. The adapt arm is \aldi; the replay arm keeps everything from the adapt arm except the adaptation.}
\label{tab:arms}
\footnotesize
\setlength{\tabcolsep}{3pt}
\begin{tabular}{@{}>{\raggedright\arraybackslash}p{0.6in}>{\raggedright\arraybackslash}p{0.95in}>{\raggedright\arraybackslash}p{0.85in}>{\raggedright\arraybackslash}p{0.75in}@{}}
\toprule
Arm & Labeled stream (2 images/step) & Unlabeled stream (2 images/step) & Updates \\
\midrule
FT & 2 labeled target images & none & max(500, 3~epochs) \\
adapt, $N = 0$ & 2 source images & full target pool, labels hidden & 2,000 \\
adapt, $N > 0$ & 1 source image + 1 of the $N$ labeled target images & full target pool & max(2,000, $2N$) \\
replay (control) & as adapt & none, no distillation & as adapt \\
\bottomrule
\end{tabular}
\end{table}

The replay arm (Table~\ref{tab:arms}) still mixes in source images and still uses strong augmentation and the EMA teacher; it drops only the unlabeled target stream and the distillation. Therefore:
\begin{itemize}
\item \textbf{adapt $-$ replay} is the value of distillation on unlabeled target images, that is, the adaptation part;
\item \textbf{replay $-$ FT} is the value of the training recipe (source-image mixing, strong augmentation, EMA).
\end{itemize}

Two parameters of \aldi were fixed by a pilot on lighting val only, under a rule written in advance: the EMA decay $\alpha = 0.9996$ and $T = 2{,}000$ updates. The run length is $\max(T, 2N)$. The original grid has 46 runs: adapt at $N = 0$ (2 seeds per shift), semi-supervised adapt at $N = 10$, 100 and 1,000 (plus 5,000 for lighting) with 3 seeds and exactly the nested subsets of FT, and replay with seed 0 at $N = 0$, 10, 100 and 1,000.

\textbf{Supplementary runs added in response to the review (30 runs):} replay with the additional seeds 1 and 2 at $N = 10$ and 100 (12 runs); adapt and replay at $N = 1{,}000$ with 3,000 updates instead of 2,000, 3 seeds (18 runs). With 3,000 updates and 1 labeled target image per step, the number of labeled target image passes equals that of FT at $N = 1{,}000$ (1,500 updates $\times$ 2 images $=$ 3,000 passes). This is the \emph{matched-passes experiment}.

\subsection{RQ5 Quantities}\label{sec:rq5quant}

\begin{itemize}
\item \textbf{\boldmath$\Delta(N)$} $=$ \aldi $-$ FT, paired by sampling seed; an advantage is meaningful when $\Delta \geq 0.005$ and the lower bound of the 95\% CI is $> 0$.
\item \textbf{Crossover \boldmath\Nx:} the smallest rung $N$ from which \aldi no longer has a meaningful advantage at any larger rung.
\item \textbf{Equivalent budget \boldmath\Neq:} the number of images FT needs to reach the same mAP as \aldi, interpolated linearly in $\log N$ at the sustained crossing of the FT curve.
\item \textbf{Replay share} $= (\text{replay} - \text{FT}) / (\text{adapt} - \text{FT})$: the part of the advantage of \aldi over FT that an arm without adaptation reproduces. It is paired by seed over the seeds common to the three arms; it is undetermined when $|\text{adapt} - \text{FT}| < 0.005$, and it is not quoted when more than 5\% of the bootstrap resamples fall into that case.
\item \textbf{Against the source model:} because FT at $N \leq 50$ is below even zero-shot, the same decomposition is recomputed with the source model as the reference, $(\text{replay} - \mAPzero) / (\text{adapt} - \mAPzero)$.
\end{itemize}

\subsection{Which Rules Were Written in Advance and Which Were Set Later}\label{sec:rules}

This distinction determines how much confidence each conclusion deserves, so it is stated separately (Table~\ref{tab:rules}).

\begin{table}[t]
\centering
\caption{Status of the analysis rules.}
\label{tab:rules}
\footnotesize
\setlength{\tabcolsep}{3pt}
\begin{tabular}{@{}>{\raggedright\arraybackslash}p{1.4in}>{\raggedright\arraybackslash}p{0.95in}>{\raggedright\arraybackslash}p{0.75in}@{}}
\toprule
Rule & When written & Nature \\
\midrule
FT schedule ($S$, $E$, LR); $\alpha$ and $T$ of \aldi & before the pilot & pre-registered \\
Definitions of $\Delta$, \Nx, \Neq; seed-0 decomposition & before the RQ5 grid & pre-registered \\
12-epoch ceiling as the denominator regardless of the result; the ``materially higher'' criterion; the reading rule for the matched-passes experiment; the rule for the main claim & 21 Nov., before the supplementary runs were scored, after the original grid had been seen & descriptive; the independent confirmation is new data \\
Log-linear fit region; interaction detection limit; replay share & 21 Nov., after the curves had been seen & descriptive \\
Recovery undetermined when $>5\%$ of bootstrap resamples are undetermined; the two ceilings presented side by side & 23 Nov., after the supplementary results had been seen & descriptive \\
Sustained crossing when computing \Neq & during the analysis of the RQ5 grid & descriptive \\
\bottomrule
\end{tabular}
\end{table}

The rule for the main claim, specifically, was written after the replay share at seed 0 had been seen (0.82--0.95 at $N = 10$ and 0.78--0.86 at $N = 100$). Its independent confirmation therefore comes from the two new seeds of the replay arm.
